\documentclass[11pt]{article}
\usepackage[margin=1in]{geometry}
\usepackage{amsmath,amssymb}
\usepackage{enumitem}
\usepackage[T1]{fontenc}

\title{COMP2300 Set 7.5 Practice Exam}
\author{Topics: Locality, Cache Basics, Associativity, Write Policies, SRAM, DRAM}
\date{}

\begin{document}
\maketitle

\noindent\textbf{Time:} 120 minutes \hfill
\textbf{Total:} 100 marks

\vspace{0.5em}
\noindent\textbf{Instructions}
\begin{itemize}[leftmargin=1.5em]
  \item Answer in English.
  \item Part A is multiple choice.
  \item Part B contains cache calculations and memory-technology reasoning.
  \item Part C contains adapted past-exam style questions with sources.
\end{itemize}

\section*{Part A: Multiple Choice (30 marks, 3 marks each)}
Choose exactly one option for each question.

\begin{enumerate}[label=\textbf{A\arabic*.}]
  \item Temporal locality means:
  \begin{enumerate}[label=(\Alph*)]
    \item Nearby addresses are likely to be used
    \item Recently used data is likely to be used again soon
    \item Memory always grows upward
    \item Every miss is compulsory
  \end{enumerate}

  \item A cache is best described as:
  \begin{enumerate}[label=(\Alph*)]
    \item A large slow memory for archival storage
    \item A small fast storage layer that keeps a subset of a larger slower memory
    \item A branch predictor
    \item A decode-stage register
  \end{enumerate}

  \item In a direct-mapped cache, each set contains:
  \begin{enumerate}[label=(\Alph*)]
    \item 0 lines
    \item 1 line
    \item 2 lines
    \item 4 lines
  \end{enumerate}

  \item Which policy updates lower memory immediately on a write hit?
  \begin{enumerate}[label=(\Alph*)]
    \item Write-back
    \item Write-through
    \item No-write-allocate
    \item Random replacement
  \end{enumerate}

  \item Which memory technology typically backs on-chip caches?
  \begin{enumerate}[label=(\Alph*)]
    \item DRAM
    \item SRAM
    \item Magnetic tape
    \item SSD flash only
  \end{enumerate}

  \item Which statement about DRAM is true?
  \begin{enumerate}[label=(\Alph*)]
    \item It stores bits in cross-coupled inverters
    \item It never needs refresh
    \item Its read operation is destructive
    \item It is always faster than SRAM
  \end{enumerate}

  \item Which component selects a row in a memory array?
  \begin{enumerate}[label=(\Alph*)]
    \item Decoder / wordline logic
    \item BTB
    \item ALU
    \item Link register
  \end{enumerate}

  \item A 2-way set associative cache allows each block to be placed in:
  \begin{enumerate}[label=(\Alph*)]
    \item Any line in the whole cache
    \item Exactly one of two lines in the indexed set
    \item Exactly one line in the whole cache
    \item Main memory only
  \end{enumerate}

  \item Which storage element is the most expensive per bit but fastest among the three?
  \begin{enumerate}[label=(\Alph*)]
    \item DRAM
    \item SRAM
    \item Flip-flops
    \item Magnetic disk
  \end{enumerate}

  \item In a memory hierarchy, the usual trend is:
  \begin{enumerate}[label=(\Alph*)]
    \item Larger memories are always faster
    \item Smaller memories tend to be faster
    \item DRAM is faster than registers
    \item Cache eliminates the need for main memory
  \end{enumerate}
\end{enumerate}

\section*{Part B: Computational Problems (50 marks)}
\begin{enumerate}[label=\textbf{B\arabic*.}]
  \item \textbf{Hit and miss rates (8 marks)}\\
  A cache sees 500 memory accesses, of which 40 are misses.
  \begin{enumerate}[label=(\alph*)]
    \item Compute the miss rate.
    \item Compute the hit rate.
  \end{enumerate}

  \item \textbf{Address breakdown (8 marks)}\\
  A direct-mapped cache has 4 sets and a block size of 8 bytes.
  \begin{enumerate}[label=(\alph*)]
    \item How many block-offset bits are required?
    \item How many set-index bits are required?
    \item In a 16-bit address, how many tag bits remain?
  \end{enumerate}

  \item \textbf{Direct-mapped placement (8 marks)}\\
  Consider the cache from B2. For address \texttt{0x0034}:
  \begin{enumerate}[label=(\alph*)]
    \item Give the binary form of the low 5 bits.
    \item Identify the block offset.
    \item Identify the set index.
  \end{enumerate}

  \item \textbf{Cache benefit calculation (8 marks)}\\
  A device has the following access costs:
  \begin{itemize}[leftmargin=1.5em]
    \item Cache hit: 1 cycle
    \item Cache miss: 4 cycles
    \item Cache disabled: 3 cycles
  \end{itemize}
  \begin{enumerate}[label=(\alph*)]
    \item Compute the average access time when hit rate is 25\%.
    \item Compute the average access time when hit rate is 50\%.
    \item For each case, state whether enabling the cache is beneficial.
  \end{enumerate}

  \item \textbf{SRAM vs DRAM reasoning (8 marks)}\\
  Explain, in one or two short sentences each:
  \begin{enumerate}[label=(\alph*)]
    \item Why SRAM is preferred for caches.
    \item Why DRAM is preferred for large main memory.
  \end{enumerate}

  \item \textbf{Write-policy comparison (10 marks)}\\
  Briefly compare the following pairs:
  \begin{enumerate}[label=(\alph*)]
    \item write-through vs write-back
    \item write-allocate vs no-write-allocate
  \end{enumerate}
  For each pair, give one advantage and one tradeoff.
\end{enumerate}

\section*{Part C: Past Exam Questions (Source-Tagged) (20 marks)}
\begin{enumerate}[label=\textbf{C\arabic*.}]
  \item \textbf{Cache enable decision (Source: exam2025-main.pdf, Q1, adapted) (8 marks)}\\
  A device access costs 3 cycles with cache disabled. With cache enabled, hits cost 1 cycle and misses cost 4 cycles. Derive the condition on hit rate under which the cache is beneficial.

  \item \textbf{Levels in the memory hierarchy (Source: 2023\_Final\_Exam\_SuppDA.pdf, microarchitecture MCQ topic, adapted) (6 marks)}\\
  List a reasonable top-to-bottom memory hierarchy for a modern processor and explain why each lower level tends to be larger but slower.

  \item \textbf{Cache misses in pipelined execution (Source: exam24r.pdf / exam24px1.pdf, Q7c context, adapted) (6 marks)}\\
  Explain why long-latency cache misses make memory-level parallelism valuable in later out-of-order designs, even though the basic cache concepts are introduced earlier.
\end{enumerate}

\end{document}
